\noindent\hspace*{\parindent}Against this we must now set the capital value of the sources of income of the widows' fund. Since the amount of the members' contributions was already deducted in the first estimate above, only the monetary assets and the pharmacy rent remain to be valued here.

The assets are stated in the last closing account, thus for 1845 July 1, as 116369 Thl. 17 Ggr. 4 Pf. But to reduce them to 1 October in accordance with the basis of calculation cited above on p. [149], we must subtract the half-yearly amount of the widows' pensions which fell due at Michaelmas 1845, 2350 Thl.; on the other hand, we must add the members' contributions for the year 1844—1845, namely 480 Thl. in full contributions and 18 Thl. 17 Ggr. 4 Pf. in partial payments, and also the half-yearly pharmacy rent which fell due at Michaelmas 1845, 500 Thl. To omit nothing, we should also have to add the interest received during the three months from 1 July to 1 October, whose amount is unknown to me; but I shall probably not be far wrong if I take one quarter of last year's interest income after deducting the collection costs, at 1163 thalers. Accordingly, omitting the fractions, I set the fund's assets on 1 October 1845 at 116171 thalers.

The pharmacy rent, at its present level of 1000 Thl., represents capital assets of 28571 Thl. or 25000 Thl., according to the assumed rate of interest. The fund's effective assets are therefore:
\begin{center}
144742 Thl. or 141171 Thl., according as one assumes \(3\tfrac{1}{2}\) or 4 per cent.
\end{center}
\noindent\hspace*{\parindent}This finally gives the fund's \textit{balance}
\begin{center}
\begin{tabular}{l}
as a \textit{deficit} of 13990 Thl. at an interest rate of \(3\tfrac{1}{2}\) per cent\\
as a \textit{surplus} of 2508 Thl. at an interest rate of 4 per cent
\end{tabular}
\end{center}
\noindent\hspace*{\parindent}It hardly needs to be remarked that the surplus arising on the second assumption is far smaller than the uncertainties remaining in the individual parts of the calculation: a single death can easily alter the entire balance at once by 3000—3500 Rthl. to the fund's disadvantage. Moreover, it must not be overlooked that I have allowed nothing at all for possible losses through the insolvency of debtors: I have omitted this because any advance estimate must remain extremely uncertain. It would be interesting to compile all the losses incurred since the foundation, but, given the form of the annual accounts, especially in earlier times, this could not be done without a very great expenditure of time: to obtain at least some idea of it, I myself have had to be content with a compilation for the last 14 years, which (if I have overlooked nothing) has given a total loss of 1729 Thl. 14 Ggr. 1 Pf., and thus an average annual loss \(=123\) Thl. 13 Ggr. Capitalized, this amounts to 3530 Thl. or 3088 Thl., according to the rate of interest; and if one were entitled to regard this average as a guide, the \textit{balance} would, at \textit{both} rates of interest, prove to be a \textit{deficit}, namely
\begin{center}
\begin{tabular}{l}
of 17520 Thl. at the interest rate of \(3\tfrac{1}{2}\) per cent\\
of\phantom{ 17} 580 Thl. at the interest rate of 4 per cent.
\end{tabular}
\end{center}
\noindent Göttingen 1845 November 2.\hfill C. F. \textsc{Gauss}.

